\begin{bt}%[1C1H5-4]%[ma:102dd419]
% nguon: Archive/SP_KhoLatex/ID6_old/1H1B5-4.tex #3352
%[HK1, Lê Quí Đôn - Đà Nẵng, 2018]%[Đức Nguyễn, 11EX1819]
Dựng về phía ngoài tam giác nhọn $ABC$ các tam giác đều $ABB'$, $ACC'$. Gọi $I$, $J$ lần lượt là trung điểm của $B'C$, $C'B$. Chứng minh rằng tam giác $AIJ$ đều.
	\loigiai{
	\begin{center}
		\begin{tikzpicture}[scale=1, font=\footnotesize,line join=round, line cap=round,>=stealth]
		\pgfmathsetmacro\a{4}
		\tkzDefPoint(0,0){B}
		\tkzDefShiftPoint[B](0:\a){C}
		\tkzDefShiftPoint[B](75:\a-0.5){A}
		\tkzDefCircle[radius](A,B)
		\tkzDefCircle[radius](B,A)
		\tkzInterCC(B,A)(A,B)    \tkzGetFirstPoint{B'}
		\tkzDefCircle[radius](A,C)
		\tkzDefCircle[radius](C,A)
		\tkzInterCC(A,C)(C,A)    \tkzGetFirstPoint{C'}
		\tkzDefMidPoint(B',C) \tkzGetPoint{I}
		\tkzDefMidPoint(B,C') \tkzGetPoint{J}
		\tkzDrawPolygon(B,B',A,C',C,B)
		\tkzDrawPolygon(A,I,J)
		\tkzDrawSegments(A,B A,C B',C B,C')
		\tkzLabelPoints[above](A,B',C',J)
		\tkzLabelPoints[below](B,C,I)
		\tkzDrawPoints[fill=black](A,B,C,B',C',I,J)
		\end{tikzpicture}
	\end{center}
Không mất tính tổng quát ta có thể giả sử tam giác $ABC$ có hướng như hình vẽ. Xét phép quay $\mathrm{Q}_{\left(A,\dfrac{\pi}{3}\right)}$ thì $B'\mapsto B$, $C\mapsto C'$, suy ra $B'C$ biến thành $BC'$ nên $I$ biến thành $J$.\\
Suy ra $AI=AJ$, $\widehat{IAJ}=\dfrac{\pi}{3}\Rightarrow $ tam giác $AIJ$ đều.
	}
\end{bt}

\begin{bt}%[1D1H2-2]%[ma:25393db1]
% nguon: Archive/SP_KhoLatex/ID6_old/1K1-1-6.tex #3708
%[DCHT Toán 11 - KNTT - Nguyễn Hoài Nam]
	Cho $\sin \alpha=-\dfrac{1}{3}$ với $\alpha \in\left(\dfrac{3 \pi}{2};2 \pi\right)$. Khi đó, giá trị của $\cos \alpha$ bằng
	\loigiai{
		Vì 	$\alpha \in\left(\dfrac{3 \pi}{2};2 \pi\right)$ nên $ \cos \alpha >0 $.\\
		Ta có $ \cos^2 \alpha=1-\sin^2 \alpha =\dfrac{8}{9}\Rightarrow \cos \alpha =\dfrac{2\sqrt{2}}{3}$.
	}
\end{bt}

\begin{bt}%[1D2V2-7]%[ma:9734a9fb]
% nguon: Archive/SP_DA_SauLuoi/1-TL-TN-THPT-VoThiSau-GHKI-NH23-24.tex #2520
%[Dự án đề kiểm tra Toán 11 GHKI NH23-24- Võ Thị Thùy Trang]%[THPT Võ Thị Sáu - Tp HCM]
	Một công ty $A$ tuyển nhân viên bằng hai phương án khi kí hợp đồng lao động.
	\begin{itemize}
		\item Phương án $1\colon $ Năm đầu tiên nhận mức lương $100$ triệu đồng, mỗi năm tiếp theo tăng thêm $12$ triệu đồng.
		\item Phương án $2\colon$ Quý đầu tiên nhận mức lương $15$ triệu đồng, mỗi quý tiếp theo tăng thêm $2{,}5$ triệu đồng.
	\end{itemize}
	Một người lao động quyết định kí hợp đồng lao động với công ty $A$ trong $10$ năm. Để tổng tiền lương nhận được trong $10$ năm tốt hơn, người lao động nên chọn phương án nào?
	\loigiai{
		\begin{itemize}
			\item Tổng tiền lương theo phương án $1$ là $S_{10}=5(2\cdot 100+9\cdot 12)=1540$ (triệu đồng).
			\item Tổng tiền lương theo phương án $2$ là $S_{40}=20(2\cdot 15+39\cdot 2{,}5 )=2550$ (triệu đồng).
		\end{itemize}
		Để tổng tiền lương nhận được trong $10$ năm tốt hơn, người lao động nên chọn phương án $2$.	
	}
\end{bt}

\begin{bt}%[1D3C2-5]%[ma:46010dd7]
% nguon: Archive/SP_SauLuoi_2024-2025/SP_HKI/1-TN-DS-TLN-THPT-ThuanThanhSo1-BacNinh-HKI-NH24-25.tex #6
%[Đề thi HK1 lớp 11 THPT Thuận Thành Số 1-24-25]%[Nguyễn Chiến]
 ($0{,}5$ điểm) Cho biết $\lim\limits_{x \rightarrow 1} \dfrac{\sqrt{ax^2+1}-bx-2}{x^3-3 x+2}$ ($a$, $b \in \mathbb{R}$) có kết quả là một số thực. Tìm $a$, $ b$.
\loigiai{
Ta có $x^3-3x+2={{\left( x-1 \right)}^2}\left( x+2 \right)$.\\
Suy ra $\lim\limits_{x \rightarrow 1} \dfrac{\sqrt{ax^2+1}-bx-2}{x^3-3 x+2}$ có kết quả là một số thực khi và chỉ khi
$\sqrt{ax^2+1}-bx-2=0$ có nghiệm kép $x=1$.\\
$\Rightarrow ax^2+1=(bx+2)^2$ có nghiệm kép $x=1$\\
$\Rightarrow \left( a-b^2 \right)x^2-4bx-3=0$ có nghiệm kép $x=1$\\
{\allowdisplaybreaks 
 \begin{eqnarray*}
&\Leftrightarrow &\left\{\begin{aligned}& a-b^2\ne 0 \\& \Delta =0 \\& a-b^2-4b-3=0\\\end{aligned} \right.
\Leftrightarrow \left\{ \begin{aligned}& a\ne b^2 \\& 16b^2+12\left(a-b^2 \right)=0 \\& a-b^2-4b-3=0 \end{aligned} \right.
\Leftrightarrow \left\{\begin{aligned}& a\ne b^2 \\& a=-\dfrac{1}{3}b^2 \\& a-b^2-4b-3=0 \end{aligned} \right.\\
&\Leftrightarrow & \left\{ \begin{aligned}& a\ne b^2 \\& a=-\dfrac{1}{3}b^2 \\& -\dfrac{1}{3}b^2-b^2-4b-3=0 \end{aligned} \right.
\Leftrightarrow \left\{ \begin{aligned}& a\ne b^2 \\& a=-\dfrac{3}{4} \\& b=-\dfrac{3}{2}\end{aligned} \right.
\Leftrightarrow \left\{ \begin{aligned}& a=-\dfrac{3}{4} \\& b=-\dfrac{3}{2}.\end{aligned} \right.
 \end{eqnarray*}  }
}
\end{bt}

\begin{bt}%[1D5H1-4]%[ma:0f465e7d]
% nguon: Archive/2024/TOAN_11_XE/TOAN_11_XE/data4/11C3/KNTT-3.OTC.tex #863
%[TeX hóa SGK KNTT]%[Phạm Hoài] 
	Người ta ghi lại tuổi thọ của một số con ong cho kết quả như sau
	\begin{center}
		\begin{tabular}{|c|c|c|c|c|c|}
			\hline 
			Tuổi thọ (ngày) & $[0;20)$ &$[20;40)$&$[40;60)$&$[60;80)$&  $[80;100)$\\
			\hline 
			Số lượng&$5$&$12$&$23$&$31$&$29$\\
			\hline 
		\end{tabular}
	\end{center}
	
	Tìm mốt của mẫu số liệu. Giải thích ý nghĩa của giá trị nhận được.
	\loigiai{
		Tần số lớn nhất là $31$ nên nhóm chứa mốt là nhóm $[60;80)$.\\
		Ta có $j=4$; $a_4=60$; $m_4=31$; $m_3=23$; $m_5=29$; $h=20$.\\
		Khi đó $M_O=60+\dfrac{31-23}{(31-23)+(31-29)}\cdot 20=76$.\\
		Như vậy, số con ong có tuổi thọ $76$ ngày là nhiều nhất.	
	} 
\end{bt}

\begin{bt}%[1D6H4-4]%[ma:2ad24017]
% nguon: Archive/SP_KhoLatex/ID6_old/KNTT6-OT-GV19.tex #3570

	Giải các phương trình sau:
	\begin{enumEX}{2}
		\item $32^{\tfrac{x+5}{x-7}}=0{,}25 \cdot 128^{\tfrac{x+17}{x-3}}$;
		\item $\log_2 x+\log_2(x-1)=1$.
	\end{enumEX}
	\loigiai{
	\begin{enumEX}{1}
		\item Điều kiện: $x \neq 3, x \neq 7$.\\
		Khi đó, ta có:
		\begin{eqnarray*}
			&&32^{\tfrac{x+5}{x-7}}=0,25 \cdot 128^{\tfrac{x+17}{x-3}}\\
			&  \Leftrightarrow& 2^{5 \cdot \tfrac{x+5}{x-7}}=2^{-2} \cdot 2^{7 \cdot \tfrac{x+17}{x-3}}\\
			&\Leftrightarrow& 2^{\tfrac{5(x+5)}{x-7}}=2^{-2+\tfrac{7(x+17)}{x-3}} \\
			& \Leftrightarrow& \dfrac{5(x+5)}{x-7}=-2+\dfrac{7(x+17)}{x-3} \\
			& \Leftrightarrow& 5(x+5)(x-3)=-2(x-7)(x-3)+7(x+17)(x-7)\\ &\Leftrightarrow& x=10.
		\end{eqnarray*}
		Kết hợp với điều kiện, ta được nghiệm của phương trình đã cho là $x=10$.
		\item Điều kiện: $x>1$.\\
		Khi đó, ta có: $\log_2 x+\log_2(x-1)=1 \Leftrightarrow \log_2 x(x-1)=1 \Leftrightarrow x(x-1)=2 \Leftrightarrow x^2-x-2=0.$\\
		Giải phương trình trên ta được hai nghiệm $x_1=-1, x_2=2$.\\
		Chỉ có nghiệm $x=2$ thoả mãn điều kiện.\\
		Vậy nghiệm của phương trình đã cho là $x=2$.
	\end{enumEX}	
}
\end{bt}

\begin{bt}%[1D7H3-1]%[ma:89675d31]
% nguon: Archive/2024/TOAN_11_XE/TOAN_11_XE/data4/11C9/KNTT-9.OTC.tex #1339
%[Tăng Xuân Phú]
	Tính đạo hàm cấp hai của các hàm số sau
	\begin{multicols}{2}
	\begin{enumerate}
		\item $y=x^3-4 x^2+2 x-3$;
		\item $y=x^2 e^x$.
	\end{enumerate}
	\end{multicols}	
	\loigiai{Đạo hàm cấp $1$ của các hàm số sau là
		\begin{enumerate}
			\item $y^\prime=3x^2-8x+2$.
			\item $y^\prime=2xe^x+x^2e^x=e^x(2x+x^2)$.
		\end{enumerate}
		Đạo hàm cấp $2$ của các hàm số trên là
		\begin{enumerate}
			\item $y^{\prime\prime}=6x-8$.
			\item $y^{\prime\prime}=e^x(2x+x^2)+e^x(2+2x)=e^x(x^2+4x+2)$.
	\end{enumerate}}
\end{bt}

\begin{bt}%[1D9V1-3]%[ma:4a6a70cd]
% nguon: Archive/2024/TOAN_10_XE/TOAN_10_XE/data4/10D9/DS10-C9-B27.tex #430
%[Hồng Trường Sơn]
    Một chiếc máy bay có hai động cơ $I$ và $II$ hoạt động độc lập với nhau. Xác suất để động cơ $I$ và  động cơ $II$ hoạt động tốt lần lượt là $0{,}8$; $0{,}7$. Tính xác suất để
\begin{enumerate}[a)]
    \item   Cả hai động cơ đều chạy tốt.
    \item Cả hai động cơ đều chạy không tốt.
\end{enumerate}     
    \loigiai{Gọi $A$, $B$ lần lượt là biến cố \lq\lq động cơ $I$, $II$ hoạt động tốt\rq\rq.
    \begin{enumerate}[a)]
        \item Gọi $C$ là biến cố \lq\lq cả hai động cơ đều chạy tốt\rq\rq.\\
        Ta có $A$ và $B$ là hai biến cố độc lập và $C=A\cdot B$.\\
         Khi đó $ \mathrm{P}(C)=\mathrm{P}(AB)=\mathrm{P}(A)\cdot \mathrm{P}(B)=0{,}8\cdot0{,}7=0{,}56$.
        \item Gọi $D$ là biến cố \lq\lq cả hai động cơ đều không chạy tốt\rq\rq.\\
        Ta có $A$ và $B$ là hai biến cố độc lập nên $\overline A$ và $\overline B$ và $C=\overline A\cdot \overline B$. \\
        Do đó $\mathrm{P}(D)=\mathrm{P}(\overline{A})\cdot \mathrm{P}\left(\overline{B}\right) =(1-0{,}8)(1-0{,}7)=0{,}06$.
    \end{enumerate} 
}
\end{bt}

\begin{bt}%[1H4V3-3]%[ma:09b4269a]
% nguon: Archive/SP_KhoLatex/ID6_old/1H2K4-3.tex #3041
%[Đề ôn số 3, Thầy Cao Thành Thái]
Cho hình chóp $S.ABCD$ có đáy $ABCD$ là hình bình hành tâm $O$. Gọi $E$, $I$, $N$ lần lượt là trung điểm của $SA$, $BC$, $CD$.
\begin{enumerate}
	\item Tìm giao tuyến của hai mặt phẳng $(ECD)$ và $(SAB)$. Suy ra giao điểm $F$ của đường thẳng $SB$ và mặt phẳng $(ECD)$.
	\item Chứng minh rằng $(OEI)\parallel (SCD)$. Tìm giao tuyến của hai mặt phẳng $(EIN)$ và $(SCD)$.
	\item Lấy điểm $H$ thuộc cạnh $SB$ sao cho $BH=2SH$, gọi $M$ là trung điểm của $AB$, $G$ là trọng tâm của tam giác $SBC$. Chứng minh rằng $AH\parallel (MNG)$.
\end{enumerate}
\loigiai
{
	\begin{enumerate}
		\item Rõ ràng $E$ là điểm chung của hai mặt phẳng $(ECD)$ và $(SAB)$. Mặt khác $CD \subset (ECD)$, $AB \subset (SAB)$ và $CD\parallel AB$ (do $ABCD$ là hình bình hành). Vậy giao tuyến của $(ECD)$ và $(SAB)$ là đường thẳng $d$ đi qua  $E$ và song song với $AB$.\\
		Trong mặt phẳng $(SAB)$ gọi $F$ là giao điểm của $d$ và $SB$. Khi đó $F\in SB$ và $F \in d$ hay $F \in SB$ và $F \in (ECD)$. Vậy $F$ là giao điểm của $SB$ và $(ECD)$.
		\begin{center}
			\begin{tikzpicture}
			\tkzDefPoints{0/0/A,6/0/D,4.5/-3/C}
			\tkzDefPointBy[translation=from D to A](C)\tkzGetPoint{B}
			\tkzDefShiftPoint[A](95:5){S}
			\tkzDefMidPoint(S,A)\tkzGetPoint{E}
			\tkzDefMidPoint(B,C)\tkzGetPoint{I}
			\tkzDefMidPoint(C,D)\tkzGetPoint{N}
			\tkzDefMidPoint(S,B)\tkzGetPoint{F}
			\tkzInterLL(A,C)(B,D)\tkzGetPoint{O}
			\tkzDefPointBy[translation=from I to E](N)\tkzGetPoint{N'}
			\tkzDefBarycentricPoint(B=1,S=2)\tkzGetPoint{H}
			\tkzDefMidPoint(A,B)\tkzGetPoint{M}
			\tkzDefMidPoint(H,B)\tkzGetPoint{K}
			\tkzCentroid(S,B,C)\tkzGetPoint{G}
			\tkzDrawPoints[fill=black](A,B,C,D,E,F,G,H,I,K,M,N,O,S)
			\tkzDrawSegments(S,B S,C S,D B,C C,D N,N' K,G S,I)
			\tkzDrawSegments[dashed](S,A A,B A,D E,C E,D O,I O,E E,I E,F M,K A,H A,C B,D M,G M,N N,G I,N E,N)
			\tkzLabelPoints[above](S)
			\tkzLabelPoints[below](B,C,I)
			\tkzLabelPoints[right](D,N)
			\tkzLabelPoints[left](A,E,H,F,K,M)
			\draw ($(O)+(0.1,-0.5)$) node{$O$};
			\draw ($(G)+(0.2,0.15)$) node{$G$};
			\end{tikzpicture}
		\end{center}
		\item Vì $OE$ là đường trung bình của tam giác $SAC$ nên $OE \parallel SC$.\\
		Vì $OI$ là đường trung bình của tam giác $BCD$ nên $OI \parallel CD$.\\
		Lại có $OE \cap OI =O$ và $SC \cap CD = C$.\\
		Vậy $(OEI)\parallel (SCD)$.\\
		Rõ ràng $N$ là điểm chung của hai mặt phẳng $(EIN)$ và $(SCD)$.\\
		Lại có $(EIN) \cap (OIE)=EI$, thêm nữa $(OIE)\parallel (SCD)$ (chứng minh trên).\\
		Vậy giao tuyến của $(EIN)$ và $(SCD)$ là đường thẳng $\Delta$ đi qua  $N$ và song song với $EI$.
		\item Gọi $K$ là trung điểm của $BH$, khi đó $SH=HK=KB$, suy ra $\dfrac{SK}{SB}=\dfrac{2}{3}$.\\
		Do $G$ là trọng tâm của tam giác $SBC$ nên $\dfrac{SG}{SI}=\dfrac{2}{3}$.\\
		Suy ra $\dfrac{SK}{SB}=\dfrac{SG}{SI}$, nên $GK\parallel BI$ hay $GK \parallel BC$.\\
		Mà $MN\parallel BC$ nên $MN \parallel BC$. Do đó $K\in (MNG)$.\\
		Ta có $AH\parallel MK$ ($MK$ là đường trung bình của tam giác $AHB$).\\
		Mà $MK \subset (MNG)$, $AH \not\subset (MNG)$ nên $AH \parallel (MNG)$.
	\end{enumerate}
	
}
\end{bt}

\begin{bt}%[1H8V2-2]%[ma:84bc8a9c]
% nguon: Archive/SP-4Phan-Dot3/Data/CD_1_OTGK2_De04.tex #4209

Cho hình chóp $S.ABCD$ có đáy là hình vuông, cạnh bên $SA$ vuông góc với mặt phẳng đáy. Gọi $H$, $K$ lần lượt là hình chiếu vuông góc của điểm $A$ lên $SC$, $SD$. Chứng minh $SC \perp (AHK)$.
\loigiai{\begin{center}
			\begin{tikzpicture}[line join=round, line cap=round,>=stealth,font=\footnotesize,scale=1]
			\def\a{3.5}
			\def\b{2.5}
			\def\h{3.5}
			\path (0:0) coordinate (A)
			++(0:\a) coordinate (D)
			++(-135:\b) coordinate (C)
			($(A)+(C)-(D)$) coordinate (B)
			($(A)!1/2!(C)$) coordinate (O)
			($(A)+(90:\h)$) coordinate (S)
			($(S)!1/3!(C)$) coordinate (H)
			($(S)!0.4!(D)$) coordinate (K);
			\draw[dashed,thick] (D)--(A) (A)--(C) (B)--(D) (A)--(B) (S)--(A) (H)--(A)--(K);
			\draw[thick] (S)--(B)--(C)--(D)--(S)--(C) (H)--(K);
			\foreach \x/\g in {A/180,B/180,C/0,D/0,O/-90,S/180,H/70,K/0}
			\fill[black] (\x) circle (1pt) ($(\g:4mm)+(\x)$) node {$\x$};	
			\newcommand{\gv}[4][black]{\draw[thick] ($(#3)!8pt!(#2)$)--($(#3)!2!($($(#3)!8pt!(#2)$)!.5!($(#3)!8pt!(#4)$)$)$)--($(#3)!8pt!(#4)$);}
			\gv{A}{H}{C}
			\gv{A}{K}{D}
		\end{tikzpicture}
\end{center}
Vì $SA \perp (ABCD)$ nên $SA \perp CD$.\\
Ta có $\heva{&CD \perp SA\\&CD \perp AD} \Rightarrow CD \perp (SAD)$.\\
Mà $AK \subset (SAD)$ nên $CD \perp AK$.\\
Ta có
$\heva{&AK \perp CD\\&AK \perp SD}\Rightarrow AK \perp (SCD)  \Rightarrow AK \perp SC$.\\
Do đó $\heva{&SC \perp AK\\&SC \perp AH} \Rightarrow SC \perp (AHK)$.
} 

\end{bt}